%% Drafted from the mapped corpus by scripts/draft_topics.py.
% Model: gpt-5.6-terra; source characters: 11878.
\begin{konspekt}
\kreq{Официално заглавие}{„Симетрични оператори в крайномерни евклидови пространства. \emph{Основни свойства.} Теорема за диагонализация.“}
\kreq{Дефиниция}{симетричен оператор --- \kr{30.2}; матрицата му спрямо ортонормиран базис --- \kref{30.2.1}.}
\kreq{Доказателство}{всички характеристични корени на симетричен оператор са реални --- \kref{30.3.1}.}
\kreq{Доказателство}{всеки два собствени вектора за различни собствени стойности са ортогонални --- \kref{30.3.2}.}
\kreq{Лема с доказателство}{инвариантност на ортогоналното допълнение --- използва се в теоремата за диагонализация --- \kref{30.4}.}
\kreq{Теорема с доказателство}{съществува ортонормиран базис, в който матрицата на симетричен оператор е диагонална --- \kref{30.5}.}
\kprereq{\kr{30.1} събира понятията от линейната алгебра --- евклидово пространство, скаларно произведение, ортонормиран базис, линеен оператор и матрицата му. Анотацията започва от \kr{30.2}.}
\kreq{Литература}{[18].}
\end{konspekt}

\section{Предварителни понятия}\label{s:30.1}

\begin{supp}[Предварителен материал]
Анотацията към Въпрос 30 започва от определението за симетричен оператор
(\S30.2). Този раздел събира понятията от линейната алгебра, които се използват
по-нататък --- евклидово пространство, скаларно произведение, ортонормиран
базис, линеен оператор и матрицата му. Той е предпоставка, а не част от
изпитния въпрос; при повторение може да се прегледа бегло.
\end{supp}

Понятията в този раздел са от курса по линейна алгебра. Изброяваме ги, защото всяко от тях се използва по-нататък в главата: без тях нито определението за симетричен оператор, нито доказателството на теоремата за диагонализация са напълно записани.

\begin{defn}[Скаларно произведение и евклидово пространство]
Нека $E$ е векторно пространство над $\mathbb{R}$. Изображението
\[
(\cdot,\cdot):E\times E\to\mathbb{R}
\]
се нарича \emph{скаларно произведение}, ако за всички $x,y,z\in E$ и $\alpha\in\mathbb{R}$ са изпълнени:

\begin{itemize}
\item симетричност: $(x,y)=(y,x)$;
\item адитивност и хомогенност по първия аргумент:
\[
(x+z,y)=(x,y)+(z,y),
\qquad
(\alpha x,y)=\alpha(x,y);
\]
\item положителна определеност: $(x,x)\geq 0$, като $(x,x)=0$ точно тогава, когато $x=0$.
\end{itemize}

От симетричността следва същото и по втория аргумент, тоест скаларното произведение е \emph{билинейно}. Двойката $\bigl(E,(\cdot,\cdot)\bigr)$ се нарича \emph{евклидово пространство}; то е крайномерно, ако $\dim E<\infty$.
\end{defn}

Стандартният пример е $\mathbb{R}^n$ със скаларно произведение
\[
(x,y)=\sum_{i=1}^n x_iy_i=x^ty,
\]
което ще наричаме \emph{стандартно евклидово скаларно произведение}.

\begin{defn}[Норма и ортогоналност]
Нека $E$ е евклидово пространство. \emph{Норма} на вектора $x$ е числото
\[
\|x\|=\sqrt{(x,x)},
\]
което е определено, защото $(x,x)\geq0$. Вектор с $\|x\|=1$ се нарича \emph{единичен}; за $x\ne0$ векторът $x/\|x\|$ е единичен и се получава чрез \emph{нормализиране} на $x$.

Векторите $x$ и $y$ са \emph{ортогонални}, което се записва $x\perp y$, ако
\[
(x,y)=0.
\]
Две подмножества $A,B\subseteq E$ са ортогонални, $A\perp B$, ако $x\perp y$ за всички $x\in A$ и $y\in B$.
\end{defn}

\begin{defn}[Ортонормиран базис]
Системата вектори $f_1,\ldots,f_k$ е \emph{ортогонална}, ако $(f_i,f_j)=0$ при $i\ne j$, и \emph{ортонормирана}, ако освен това всички вектори са единични, тоест
\[
(f_i,f_j)=\delta_{ij}=
\begin{cases}
1, & i=j,\\
0, & i\ne j.
\end{cases}
\]
\emph{Ортонормиран базис} на $E$ е базис, който е ортонормирана система.
\end{defn}

Ортонормираният базис е удобен именно защото свежда скаларното произведение до координатно пресмятане: ако $e_1,\ldots,e_n$ е ортонормиран базис и
\[
x=\sum_{i=1}^n x_ie_i,
\qquad
y=\sum_{j=1}^n y_je_j,
\]
то от билинейността
\[
(x,y)=\sum_{i=1}^n\sum_{j=1}^n x_iy_j(e_i,e_j)=\sum_{i=1}^n x_iy_i=x^ty .
\]
В частност $x_i=(x,e_i)$: координатите на вектор спрямо ортонормиран базис са скаларните му произведения с базисните вектори.

\begin{prop}[Съществуване на ортонормиран базис]
Всяко крайномерно евклидово пространство $E$ с $\dim E=n\geq1$ притежава ортонормиран базис.
\end{prop}

\begin{proof}
Нека $v_1,\ldots,v_n$ е произволен базис на $E$. По метода на Грам--Шмит строим индуктивно
\[
w_1=v_1,
\qquad
w_k=v_k-\sum_{i=1}^{k-1}\frac{(v_k,w_i)}{(w_i,w_i)}\,w_i,
\qquad k=2,\ldots,n.
\]
Допускаме индуктивно, че $w_1,\ldots,w_{k-1}$ са ненулеви, попарно ортогонални и
\[
\mathcal{L}(w_1,\ldots,w_{k-1})=\mathcal{L}(v_1,\ldots,v_{k-1}).
\]
Тогава дробите са определени, защото $(w_i,w_i)>0$ за ненулево $w_i$. За $j<k$ пресмятаме
\[
(w_k,w_j)
=(v_k,w_j)-\sum_{i=1}^{k-1}\frac{(v_k,w_i)}{(w_i,w_i)}(w_i,w_j)
=(v_k,w_j)-(v_k,w_j)=0,
\]
тъй като в сумата оцелява само събираемото с $i=j$. Освен това $w_k\ne0$: иначе $v_k$ би лежал в $\mathcal{L}(w_1,\ldots,w_{k-1})=\mathcal{L}(v_1,\ldots,v_{k-1})$, което противоречи на линейната независимост на базиса. От построението следва и равенството на обвивките за $k$.

Накрая полагаме $e_i=w_i/\|w_i\|$. Системата $e_1,\ldots,e_n$ е ортонормирана. Всяка ортонормирана система е линейно независима: ако $\sum_i\alpha_ie_i=0$, скаларното умножение с $e_j$ дава $\alpha_j=0$. Значи $n$ такива вектора образуват базис на $E$.
\end{proof}

\begin{defn}[Линейна обвивка и ортогонално допълнение]
С $\mathcal{L}(v_1,\ldots,v_k)$ означаваме \emph{линейната обвивка} на векторите $v_1,\ldots,v_k$, тоест множеството от всички техни линейни комбинации; това е най-малкото подпространство, което ги съдържа. В частност за $u\ne0$
\[
\mathcal{L}(u)=\{\alpha u\mid\alpha\in\mathbb{R}\}
\]
е правата, породена от $u$.

\emph{Ортогонално допълнение} на подпространството $U\subseteq E$ е множеството
\[
U^\perp=\{x\in E\mid (x,u)=0\ \text{за всяко }u\in U\}.
\]
То е подпространство на $E$, защото условието е линейно по $x$.
\end{defn}

\begin{prop}[Ортогонално разлагане]
Нека $E$ е крайномерно евклидово пространство и $U\subseteq E$ е подпространство. Тогава
\[
E=U\oplus U^\perp,
\qquad
\dim U^\perp=\dim E-\dim U.
\]
\end{prop}

\begin{proof}
Ако $U=\{0\}$, твърдението е очевидно, защото тогава $U^\perp=E$. Нека $\dim U=k\geq1$ и нека $u_1,\ldots,u_k$ е ортонормиран базис на $U$; такъв съществува по предходното твърдение, приложено към $U$ със същото скаларно произведение.

За произволно $x\in E$ полагаме
\[
p=\sum_{i=1}^k (x,u_i)u_i\in U,
\qquad
r=x-p.
\]
За всяко $j$ имаме
\[
(r,u_j)=(x,u_j)-\sum_{i=1}^k (x,u_i)(u_i,u_j)=(x,u_j)-(x,u_j)=0,
\]
защото базисът е ортонормиран. Понеже $r$ е ортогонален на всички $u_j$, по линейност той е ортогонален на всяка тяхна линейна комбинация, тоест $r\in U^\perp$. Следователно $x=p+r$ и
\[
E=U+U^\perp.
\]

Сумата е директна: ако $x\in U\cap U^\perp$, то $x$ е ортогонален на себе си, тоест $(x,x)=0$, откъдето $x=0$ поради положителната определеност. Оттук следва и равенството за размерностите.
\end{proof}

\begin{defn}[Собствена стойност, собствен вектор, собствено подпространство]
Нека $\varphi:E\to E$ е линеен оператор. Числото $\lambda\in\mathbb{R}$ е \emph{собствена стойност} на $\varphi$, ако съществува вектор $x\ne0$, за който
\[
\varphi(x)=\lambda x.
\]
Всеки такъв $x$ се нарича \emph{собствен вектор} за $\lambda$. Множеството
\[
E_\lambda=\{x\in E\mid \varphi(x)=\lambda x\}
\]
е подпространство на $E$ и се нарича \emph{собствено подпространство} за $\lambda$; то се състои от собствените вектори за $\lambda$ и нулевия вектор.
\end{defn}

\begin{defn}[Характеристичен полином и характеристичен корен]
Нека $A\in M_{n\times n}(\mathbb{R})$. Полиномът
\[
\chi_A(\lambda)=\det(A-\lambda I)
\]
от степен $n$ се нарича \emph{характеристичен полином} на $A$. Неговите корени в $\mathbb{C}$ се наричат \emph{характеристични корени}.

Подобните матрици имат един и същ характеристичен полином, затова матриците на един и същ оператор спрямо различни базиси дават един и същ полином. Той се нарича характеристичен полином на оператора $\varphi$.
\end{defn}

Връзката между двете понятия е следната. Ако $A$ е матрицата на $\varphi$ спрямо някакъв базис, то за $\lambda\in\mathbb{R}$
\[
\varphi(x)=\lambda x\ \text{за някое }x\ne0
\iff
(A-\lambda I)X=0\ \text{има ненулево решение}
\iff
\det(A-\lambda I)=0 .
\]
Следователно \emph{собствените стойности на $\varphi$ са точно реалните характеристични корени}. Затова твърдението, че всички характеристични корени на симетричен оператор са реални, е по същество твърдение за пълнотата на спектъра му.

\begin{remark}
По основната теорема на алгебрата всеки полином от степен $n\geq1$ с комплексни коефициенти има корен в $\mathbb{C}$. Затова всеки оператор в пространство с $\dim E\geq1$ има поне един характеристичен корен, но той не е задължително реален; именно това е препятствието, което симетричността премахва.
\end{remark}

\begin{defn}[Инвариантно подпространство и ограничение]
Подпространството $U\subseteq E$ е \emph{инвариантно} относно линейния оператор $\varphi:E\to E$, ако
\[
\varphi(U)\subseteq U.
\]
Тогава \emph{ограничението}
\[
\varphi|_U:U\to U,
\qquad
\varphi|_U(x)=\varphi(x),
\]
е коректно определен линеен оператор в $U$.
\end{defn}

\begin{defn}[Ортогонална матрица]
Матрицата $Q\in M_{n\times n}(\mathbb{R})$ се нарича \emph{ортогонална}, ако
\[
Q^tQ=I,
\]
тоест $Q^{-1}=Q^t$. Еквивалентно, стълбовете на $Q$ образуват ортонормирана система спрямо стандартното скаларно произведение в $\mathbb{R}^n$.
\end{defn}

\section{Симетрични оператори}\label{s:30.2}

Оттук нататък $E$ е крайномерно евклидово пространство над $\mathbb{R}$ със скаларно произведение $(\cdot,\cdot)$ и $\dim E=n$.

\begin{defn}
Линейният оператор $\varphi:E\to E$ се нарича \emph{симетричен}, ако за всеки два вектора $x,y\in E$ е изпълнено
\[
(\varphi(x),y)=(x,\varphi(y)).
\]
\end{defn}

Това определение изразява съгласуваност между действието на оператора и скаларното произведение. При симетричен оператор прехвърлянето на оператора от първия към втория аргумент на скаларното произведение не променя стойността му.

\begin{remark}
В по-общия случай на унитарно пространство над $\mathbb{C}$ аналогичното понятие е \emph{ермитов оператор}. Тогава условието е
\[
\langle\varphi(x),y\rangle=\langle x,\varphi(y)\rangle.
\]
В настоящата глава разглеждаме евклидови пространства над $\mathbb{R}$, поради което използваме термина \emph{симетричен оператор}.
\end{remark}

\subsection{Матрица на симетричен оператор}\label{s:30.2.1}

Нека $e_1,\ldots,e_n$ е ортонормиран базис на $E$ и
\[
A=(a_{ij})\in M_{n\times n}(\mathbb{R})
\]
е матрицата на $\varphi$ спрямо този базис. По определение координатите на $\varphi(e_j)$ са елементите от $j$-тия стълб на $A$, тоест
\[
\varphi(e_j)=\sum_{i=1}^n a_{ij}e_i.
\]

\begin{defn}
Матрицата $A=(a_{ij})\in M_{n\times n}(\mathbb{R})$ се нарича \emph{симетрична}, ако
\[
A^t=A.
\]
Еквивалентно,
\[
a_{ij}=a_{ji}\qquad (1\leq i,j\leq n).
\]
\end{defn}

\begin{prop}
Нека $\varphi:E\to E$ е линеен оператор и $e_1,\ldots,e_n$ е ортонормиран базис на $E$. Тогава следните условия са еквивалентни:

\begin{enumerate}
\item $\varphi$ е симетричен оператор;
\item матрицата на $\varphi$ спрямо базиса $e_1,\ldots,e_n$ е симетрична.
\end{enumerate}
\end{prop}

\begin{proof}
Нека първо $\varphi$ е симетричен оператор и $A=(a_{ij})$ е неговата матрица спрямо дадения ортонормиран базис. За произволни $i,j\in\{1,\ldots,n\}$ имаме
\[
a_{ji}=(\varphi(e_i),e_j).
\]
Тъй като $\varphi$ е симетричен,
\[
(\varphi(e_i),e_j)=(e_i,\varphi(e_j)).
\]
От ортонормираността на базиса следва
\[
(e_i,\varphi(e_j))=a_{ij}.
\]
Следователно
\[
a_{ji}=a_{ij}
\]
за всички $i,j$, откъдето $A^t=A$.

Обратно, нека матрицата $A$ е симетрична. Тогава
\[
(\varphi(e_i),e_j)=a_{ji}=a_{ij}=(e_i,\varphi(e_j))
\]
за всички базисни вектори $e_i,e_j$.

Нека сега
\[
x=\sum_{i=1}^n x_i e_i,\qquad
y=\sum_{j=1}^n y_j e_j.
\]
Използвайки линейността на $\varphi$ и билинейността на скаларното произведение, получаваме
\[
\begin{aligned}
(\varphi(x),y)
&=
\left(
\varphi\left(\sum_{i=1}^n x_i e_i\right),
\sum_{j=1}^n y_j e_j
\right) \\
&=
\sum_{i=1}^n\sum_{j=1}^n
x_i y_j(\varphi(e_i),e_j) \\
&=
\sum_{i=1}^n\sum_{j=1}^n
x_i y_j(e_i,\varphi(e_j)) \\
&=
\left(
\sum_{i=1}^n x_i e_i,
\varphi\left(\sum_{j=1}^n y_j e_j\right)
\right) \\
&=(x,\varphi(y)).
\end{aligned}
\]
Следователно $\varphi$ е симетричен оператор.
\end{proof}

\begin{cor}
Линеен оператор в крайномерно евклидово пространство е симетричен точно тогава, когато матрицата му спрямо произволен ортонормиран базис е симетрична.
\end{cor}

\begin{proof}
В ортонормиран базис скаларното произведение има вид $x^ty$. Ако $A$ е матрицата на оператора, условието за симетричност е $(Ax)^ty=x^tAy$ за всички $x,y$, тоест $x^t(A^t-A)y=0$. Избирането на координатни вектори за $x,y$ показва, че това е еквивалентно на $A^t=A$. Аргументът важи във всеки ортонормиран базис.
\end{proof}

\begin{supp}[Бележка]
За матрици с комплексни елементи условието
\[
\overline{A}^{\,t}=A
\]
описва ермитова, а не непременно симетрична матрица. В реалния случай $\overline{A}=A$, така че това условие се свежда до $A^t=A$.
\end{supp}

\section{Собствени стойности и собствени вектори}\label{s:30.3}

Нека $\varphi:E\to E$ е симетричен оператор. Основните му спектрални свойства са, че характеристичните му корени са реални и че собствените подпространства за различни собствени стойности са ортогонални.

\subsection{Реалност на характеристичните корени}\label{s:30.3.1}

\begin{thm}
Всички характеристични корени на симетричен оператор в крайномерно евклидово пространство са реални числа.
\end{thm}

\begin{proof}
Нека $A$ е матрицата на $\varphi$ спрямо ортонормиран базис на $E$. Съгласно предходното твърдение $A$ е реална симетрична матрица.

Нека $\lambda\in\mathbb{C}$ е характеристичен корен на $A$. Тогава съществува ненулев комплексен стълб-вектор
\[
z=
\begin{pmatrix}
z_1\\
\vdots\\
z_n
\end{pmatrix}
\in\mathbb{C}^n
\]
такъв, че
\[
Az=\lambda z.
\]
Умножаваме отляво с ред-вектора
\[
\overline{z}^{\,t}=
(\overline{z_1},\ldots,\overline{z_n}).
\]
Получаваме
\[
\overline{z}^{\,t}Az
=
\lambda\,\overline{z}^{\,t}z
=
\lambda\sum_{i=1}^n |z_i|^2.
\]
Нека
\[
q=\overline{z}^{\,t}Az.
\]
Тъй като $A$ е реална симетрична матрица и $q$ е скалар,
\[
\overline q
=z^tA\overline z
=(z^tA\overline z)^t
=\overline z^{\,t}A^tz
=\overline z^{\,t}Az
=q.
\]
Следователно $q$ е реално число. От друга страна, комплексното спрягане на равенството $Az=\lambda z$ дава
\[
A\overline{z}=\overline{\lambda}\,\overline{z}.
\]
Следователно
\[
q=\overline q
=z^tA\overline z
=
\overline{\lambda}\sum_{i=1}^n |z_i|^2.
\]
Значи
\[
\lambda\sum_{i=1}^n |z_i|^2
=
\overline{\lambda}\sum_{i=1}^n |z_i|^2.
\]
Понеже $z\ne0$, то
\[
\sum_{i=1}^n |z_i|^2>0.
\]
Следователно
\[
\lambda=\overline{\lambda},
\]
което означава, че $\lambda\in\mathbb{R}$.
\end{proof}

\begin{remark}
Доказателството разглежда характеристичните корени като комплексни числа. Това е необходимо, защото характеристичният полином на реална матрица може да има комплексни корени. За симетрична матрица доказаното свойство показва, че всички тези корени всъщност са реални.
\end{remark}

\begin{cor}
Всеки симетричен оператор $\varphi:E\to E$ има собствена стойност, ако $\dim E\geq1$.
\end{cor}

\begin{proof}
Характеристичният полином на матрицата на $\varphi$ има комплексен корен. По теоремата всички негови корени са реални, следователно този корен е реална собствена стойност на $\varphi$.
\end{proof}

\subsection{Ортогоналност на собствени вектори}\label{s:30.3.2}

\begin{prop}
Нека $\varphi:E\to E$ е симетричен оператор. Ако $u,v\in E\setminus\{0\}$ са собствени вектори, съответстващи съответно на различните собствени стойности $\lambda$ и $\mu$, то
\[
u\perp v.
\]
\end{prop}

\begin{proof}
Имаме
\[
\varphi(u)=\lambda u,\qquad
\varphi(v)=\mu v,
\]
където $\lambda\ne\mu$. От симетричността на $\varphi$ следва
\[
(\varphi(u),v)=(u,\varphi(v)).
\]
След заместване получаваме
\[
(\lambda u,v)=(u,\mu v),
\]
тоест
\[
\lambda(u,v)=\mu(u,v).
\]
Следователно
\[
(\lambda-\mu)(u,v)=0.
\]
Понеже $\lambda\ne\mu$, имаме
\[
(u,v)=0.
\]
Следователно $u\perp v$.
\end{proof}

\begin{cor}
Собствените подпространства на симетричен оператор, съответстващи на различни собствени стойности, са ортогонални.
\end{cor}

\begin{proof}
Нека $E_\lambda$ и $E_\mu$ са собствените подпространства за $\lambda$ и $\mu$, където $\lambda\ne\mu$. Всеки ненулев вектор от $E_\lambda$ е собствен вектор за $\lambda$, а всеки ненулев вектор от $E_\mu$ е собствен вектор за $\mu$. По предходното твърдение всеки такъв чифт вектори е ортогонален. Следователно
\[
E_\lambda\perp E_\mu.
\]
\end{proof}

\begin{remark}
Условието за различност на собствените стойности е съществено. Два собствени вектора, отговарящи на една и съща собствена стойност, не са задължително ортогонални. В собственото подпространство обаче може да се избере ортонормиран базис.
\end{remark}

\section{Инвариантност на ортогоналното допълнение}\label{s:30.4}

За доказателството на теоремата за диагонализация е необходимо свойството, че ортогоналното допълнение на собствен вектор е инвариантно относно симетричния оператор.

\begin{lem}
Нека $\varphi:E\to E$ е симетричен оператор и $u\in E\setminus\{0\}$ е собствен вектор:
\[
\varphi(u)=\lambda u.
\]
Тогава ортогоналното допълнение
\[
\mathcal{L}(u)^\perp
\]
е инвариантно относно $\varphi$, тоест
\[
\varphi\bigl(\mathcal{L}(u)^\perp\bigr)
\subseteq
\mathcal{L}(u)^\perp.
\]
\end{lem}

\begin{proof}
Нека $x\in\mathcal{L}(u)^\perp$. Тогава
\[
(x,u)=0.
\]
Ще докажем, че $\varphi(x)\in\mathcal{L}(u)^\perp$. От симетричността на $\varphi$ следва
\[
(\varphi(x),u)=(x,\varphi(u)).
\]
Тъй като $\varphi(u)=\lambda u$, получаваме
\[
(\varphi(x),u)=(x,\lambda u)=\lambda(x,u)=0.
\]
Следователно
\[
\varphi(x)\perp u,
\]
тоест
\[
\varphi(x)\in\mathcal{L}(u)^\perp.
\]
\end{proof}

\begin{lem}
При условията на предходната лема ограничението
\[
\varphi\big|_{\mathcal{L}(u)^\perp}:
\mathcal{L}(u)^\perp\to\mathcal{L}(u)^\perp
\]
е симетричен оператор.
\end{lem}

\begin{proof}
Нека $x,y\in\mathcal{L}(u)^\perp$. Понеже $\varphi$ е симетричен в $E$, имаме
\[
(\varphi(x),y)=(x,\varphi(y)).
\]
Съгласно предходната лема $\varphi(x)$ и $\varphi(y)$ принадлежат на $\mathcal{L}(u)^\perp$, следователно това равенство е именно условието ограничението на $\varphi$ върху това подпространство да е симетрично.
\end{proof}

\section{Теорема за диагонализация}\label{s:30.5}

\begin{keythm}[Теорема за диагонализация на симетричен оператор]
Нека $E$ е крайномерно евклидово пространство и $\varphi:E\to E$ е симетричен оператор. Тогава съществува ортонормиран базис
\[
e_1,\ldots,e_n
\]
на $E$, съставен от собствени вектори на $\varphi$.

Еквивалентно, съществува ортонормиран базис на $E$, спрямо който матрицата на $\varphi$ е диагонална:
\[
[\varphi]_{(e_1,\ldots,e_n)}
=
\begin{pmatrix}
\lambda_1 & 0 & \cdots & 0\\
0 & \lambda_2 & \cdots & 0\\
\vdots & \vdots & \ddots & \vdots\\
0 & 0 & \cdots & \lambda_n
\end{pmatrix},
\]
където $\lambda_1,\ldots,\lambda_n\in\mathbb{R}$ са собствени стойности на $\varphi$, като някои от тях могат да съвпадат.
\end{keythm}

\begin{proof}
Ще докажем твърдението с индукция по
\[
n=\dim E.
\]

При $n=1$ твърдението е очевидно. Ако $e_1$ е единичен вектор, то $E=\mathcal{L}(e_1)$ и поради линейността на $\varphi$ съществува $\lambda_1\in\mathbb{R}$, за което
\[
\varphi(e_1)=\lambda_1e_1.
\]
Следователно $e_1$ е ортонормиран базис от собствени вектори.

Нека твърдението е вярно за всички евклидови пространства с размерност $n-1$, и нека сега $\dim E=n$.

По доказаното по-горе симетричният оператор $\varphi$ има реална собствена стойност $\lambda_1$. Нека
\[
u_1\in E\setminus\{0\}
\]
е съответен собствен вектор:
\[
\varphi(u_1)=\lambda_1u_1.
\]
Нормализираме $u_1$ и полагаме
\[
e_1=\frac{u_1}{\|u_1\|}.
\]
Тогава $\|e_1\|=1$ и
\[
\varphi(e_1)=\lambda_1e_1.
\]

Нека
\[
U=\mathcal{L}(e_1)^\perp.
\]
По твърдението за ортогоналното разлагане $E=\mathcal{L}(e_1)\oplus U$, а тъй като $\dim\mathcal{L}(e_1)=1$, това е подпространство с размерност $n-1$. От лемата следва, че $U$ е инвариантно относно $\varphi$. Следователно ограничението
\[
\varphi|_U:U\to U
\]
е коректно определен линеен оператор. Освен това то е симетрично.

По индукционното предположение за оператора $\varphi|_U$ съществува ортонормиран базис
\[
e_2,\ldots,e_n
\]
на $U$, съставен от собствени вектори на ограничението. Следователно съществуват реални числа $\lambda_2,\ldots,\lambda_n$, за които
\[
\varphi(e_i)=\lambda_i e_i,
\qquad i=2,\ldots,n.
\]

Тъй като разлагането
\[
E=\mathcal{L}(e_1)\oplus U
\]
е ортогонално --- всеки вектор от $U$ е ортогонален на $e_1$ по определението за $U$ --- системата
\[
e_1,e_2,\ldots,e_n
\]
е ортонормиран базис на $E$. Всеки негов вектор е собствен вектор на $\varphi$. Следователно матрицата на $\varphi$ спрямо този базис е
\[
\begin{pmatrix}
\lambda_1 & 0 & \cdots & 0\\
0 & \lambda_2 & \cdots & 0\\
\vdots & \vdots & \ddots & \vdots\\
0 & 0 & \cdots & \lambda_n
\end{pmatrix}.
\]
Така индукционната стъпка е доказана.
\end{proof}

\begin{cor}
Всеки симетричен оператор в крайномерно евклидово пространство е диагонализируем спрямо ортонормиран базис.
\end{cor}

\begin{proof}
По теоремата за диагонализация съществува ортонормиран базис от собствени вектори $e_1,\ldots,e_n$. Ако $\varphi(e_i)=\lambda_i e_i$, $i$-тата колона на матрицата в този базис има единствена възможна ненулева координата $\lambda_i$ на ред $i$. Следователно матрицата е $\operatorname{diag}(\lambda_1,\ldots,\lambda_n)$.
\end{proof}

\begin{remark}
Теоремата не твърди само съществуване на базис, в който матрицата е диагонална. Същественото допълнително свойство е, че базисът може да бъде избран ортонормиран. Именно това е специфичната особеност на симетричните оператори.
\end{remark}

\subsection{Връзка със симетричните матрици}\label{s:30.5.1}

Нека $A\in M_{n\times n}(\mathbb{R})$ е симетрична матрица. Разглеждаме линейния оператор
\[
\varphi:\mathbb{R}^n\to\mathbb{R}^n,
\qquad
\varphi(x)=Ax,
\]
спрямо стандартното евклидово скаларно произведение. Матрицата на $\varphi$ спрямо каноничния ортонормиран базис е $A$, следователно $\varphi$ е симетричен оператор.

От теоремата за диагонализация следва, че съществува ортонормиран базис от собствени вектори на $A$. Ако $Q$ е матрицата, чиито стълбове са тези базисни вектори, то $Q$ е ортогонална матрица и в този базис матрицата на оператора е диагонална. Следователно симетричната матрица може да бъде приведена до диагонален вид чрез преход към ортонормиран базис.

\begin{example}
Нека
\[
A=
\begin{pmatrix}
2 & 1\\
1 & 2
\end{pmatrix}.
\]
Матрицата е симетрична, защото
\[
A^t=
\begin{pmatrix}
2 & 1\\
1 & 2
\end{pmatrix}
=A.
\]
Характеристичният полином е
\[
\det(A-\lambda I)
=
\begin{vmatrix}
2-\lambda & 1\\
1 & 2-\lambda
\end{vmatrix}
=
(2-\lambda)^2-1.
\]
Следователно характеристичните корени са
\[
\lambda_1=3,\qquad \lambda_2=1.
\]

За $\lambda_1=3$ получаваме собствен вектор
\[
u_1=
\begin{pmatrix}
1\\
1
\end{pmatrix},
\]
а за $\lambda_2=1$ --- собствен вектор
\[
u_2=
\begin{pmatrix}
1\\
-1
\end{pmatrix}.
\]
Техният скаларен продукт е
\[
(u_1,u_2)=1\cdot1+1\cdot(-1)=0,
\]
което илюстрира ортогоналността на собствени вектори за различни собствени стойности.

След нормализиране получаваме ортонормирани собствени вектори
\[
e_1=\frac{1}{\sqrt{2}}
\begin{pmatrix}
1\\
1
\end{pmatrix},
\qquad
e_2=\frac{1}{\sqrt{2}}
\begin{pmatrix}
1\\
-1
\end{pmatrix}.
\]
Спрямо ортонормирания базис $e_1,e_2$ матрицата на оператора е
\[
\begin{pmatrix}
3 & 0\\
0 & 1
\end{pmatrix}.
\]
\end{example}

\section{Обобщение на свойствата}\label{s:30.6}

Нека $\varphi:E\to E$ е симетричен оператор в крайномерно евклидово пространство. Тогава са изпълнени следните основни свойства:

\begin{enumerate}
\item За всеки $x,y\in E$ е изпълнено
\[
(\varphi(x),y)=(x,\varphi(y)).
\]

\item Матрицата на $\varphi$ спрямо всеки ортонормиран базис е реална симетрична матрица.

\item Всички характеристични корени, съответно всички собствени стойности, на $\varphi$ са реални.

\item Ако $u$ и $v$ са собствени вектори за различни собствени стойности, то
\[
(u,v)=0.
\]

\item Ако $u$ е собствен вектор, то подпространството
\[
\mathcal{L}(u)^\perp
\]
е инвариантно относно $\varphi$.

\item Съществува ортонормиран базис на $E$, съставен от собствени вектори на $\varphi$.

\item Спрямо този ортонормиран базис матрицата на $\varphi$ е диагонална, а елементите по главния диагонал са реалните собствени стойности на оператора.
\end{enumerate}

\section{Изпитен фокус}\label{s:30.7}

\begin{itemize}
\item Да се даде определението за симетричен оператор:
\[
(\varphi(x),y)=(x,\varphi(y))
\qquad
\text{за всички }x,y\in E.
\]

\item Да се знае, че матрицата на симетричен оператор спрямо ортонормиран базис е симетрична, и да се възпроизведе доказателството чрез равенството
\[
a_{ji}=(\varphi(e_i),e_j)=(e_i,\varphi(e_j))=a_{ij}.
\]

\item Да се знае и обратното: симетрична матрица спрямо ортонормиран базис определя симетричен оператор.

\item Да се формулира и докаже, че всички характеристични корени на симетричен оператор са реални, като се използва равенството
\[
\lambda\sum_{i=1}^n|z_i|^2
=
\overline{\lambda}\sum_{i=1}^n|z_i|^2.
\]

\item Да се формулира и докаже ортогоналността на собствени вектори за различни собствени стойности:
\[
(\lambda-\mu)(u,v)=0.
\]

\item Да се знае, че ако $u$ е собствен вектор, то $\mathcal{L}(u)^\perp$ е инвариантно относно симетричния оператор, защото
\[
(\varphi(x),u)=(x,\varphi(u))=\lambda(x,u)=0.
\]

\item Да се формулира теоремата за диагонализация: всеки симетричен оператор в крайномерно евклидово пространство има ортонормиран базис от собствени вектори.

\item Да се възпроизведе доказателствената идея за теоремата: индукция по $\dim E$, избор и нормализиране на собствен вектор, преминаване към неговото ортогонално допълнение и прилагане на индукционното предположение върху ограничението на оператора.
\end{itemize}
